\section{外部接口与界面需求}

\subsection{四端页面结构}

\srsfigure{figures/srs/15-page-information-architecture.pdf}{四端页面信息架构图}{fig:page-ia}

顾客端以店铺浏览、购物车、订单和个人中心为主导航；商家端以店铺、商品、订单和看板为主；骑手端以可接任务、当前任务和历史为主；管理端以审核、治理和平台看板为主。角色切换后必须进入相应工作台，不保留上一角色的私有页面状态。

\subsection{关键页面原型}

\srsfigure{figures/srs/16-key-page-wireframes.pdf}{关键页面低保真原型组}{fig:key-wireframes}

原型组至少包含顾客 AI 点餐页、订单履约详情页、商家订单工作台、骑手当前任务页和管理员审核/看板页。原型用于约定信息层级和状态反馈，视觉色彩可在不改变需求的前提下调整。

\subsection{页面通用状态}

\begin{longtable}{|L{0.17\linewidth}|L{0.31\linewidth}|L{0.42\linewidth}|}
\caption{页面通用状态}\label{tab:ui-state}\\ \hline
\textbf{状态} & \textbf{展示要求} & \textbf{交互要求} \\ \hline
加载中 & 显示骨架屏或明确加载指示 & 防止重复提交产生副作用的操作 \\ \hline
空数据 & 说明当前无数据的原因 & 在合适场景提供创建、搜索或返回入口 \\ \hline
表单错误 & 定位到具体字段并显示可理解文字 & 保留其他合法输入，不清空整张表单 \\ \hline
权限不足 & 显示无权访问，不展示敏感数据 & 引导返回当前角色首页 \\ \hline
状态冲突 & 显示数据已变更或操作已被他人完成 & 刷新最新状态，不覆盖已成功结果 \\ \hline
服务失败 & 显示简洁错误与重试入口 & 写操作仅在存在幂等保护时允许自动重试 \\ \hline
\end{longtable}

\subsection{REST 接口约束}

\begin{longtable}{|L{0.22\linewidth}|L{0.68\linewidth}|}
\caption{REST 接口约束}\label{tab:rest-constraints}\\ \hline
\textbf{项目} & \textbf{要求} \\ \hline
路径 & 新接口使用 \texttt{/api/v1} 版本前缀与资源名词；旧 \texttt{/api} 接口在迁移期保持兼容，不要求一次性改完全部路径。 \\ \hline
HTTP 方法 & GET 查询，POST 创建，PATCH 局部修改，DELETE 删除/逻辑删除；支付、取消、接单等动作可使用动作子资源。 \\ \hline
状态码 & 200/201 表示成功，400 表示参数错误，401 表示未认证，403 表示无权，404 表示资源不存在，409 表示状态/并发冲突，500 表示未处理服务端错误。 \\ \hline
统一响应 & 包含业务码、消息、数据和请求标识；不向客户端暴露堆栈和数据库错误。 \\ \hline
分页 & 页码从 1 开始，默认每页 20 条，单页最多 100 条；响应包含总数、页码、页大小和列表。 \\ \hline
认证与归属 & 需要登录的接口从认证上下文获取用户身份，不信任客户端传入的用户 ID；角色校验后继续校验资源归属。 \\ \hline
\end{longtable}

\subsection{外部服务与降级}

\srsfigure{figures/srs/17-external-services.pdf}{外部服务依赖与降级图}{fig:external-services}

AI、图片识别、语音识别、地图和对象存储必须由后端适配层封装。自动化测试使用可控的 Mock/Stub；任一外部服务失败时，不得导致登录、普通搜索、购物车、订单或文本地址功能不可用。
